%% ============================================================
%% MODELO DE ATIVIDADE · GESTÃO E LIDERANÇA
%%
%% Copie este arquivo para atividades/ (atividade-02.tex, por exemplo),
%% inclua-o no relatorio-aacc.tex com \input{atividades/atividade-02}
%% e escreva. Cada \preencher{…} é uma pendência: troque-o pelo seu
%% texto. Cada \figuraprovisoria também: troque-a pela figura de verdade.
%%
%% Uma atividade é UMA entrega sua, com começo, meio e fim. "Participei
%% do projeto X por dois anos" não é uma atividade: são várias. Divida.
%% ============================================================
\begin{atividade}{Supervisão do projeto P (ou gerência do departamento D)}{
  tipo         = gestao,
  modalidade   = {},   % a chave da modalidade pedida: EXT, PES, COMP, VPC, EVT, PRD, ORG, ENS, GES, CUR, OUT
  alternativa  = {},   % a modalidade alternativa, se o Colegiado não aceitar a primeira
  horas        = {},   % as horas desta atividade, com base nos registros (Parte IV)
  inicio       = {},   % AAAA-MM
  fim          = {},   % AAAA-MM, ou: em andamento
  projeto      = {},
  papel        = {},   % Responsável, Corresponsável ou Colaborador(a)
  supervisor   = {},   % quem confirma: nome e cargo (o supervisor do projeto, o gerente)
  registros    = {},   % os códigos no SOMA, com ponto e vírgula: NBL-14; NRO-PRO-003-7
  competencias = {},   % as competências das DCN que ela desenvolveu (I a VIII): I, III, V
  disciplinas  = {},   % as disciplinas do seu curso ligadas a ela: ELE075 Eletrônica; ELE067 Circuitos
}

\begin{contexto}
\preencher{O cargo, o time e o momento da equipe quando você assumiu.}
\end{contexto}

\begin{contribuicao}
\preencher{As suas decisões e entregas como líder.}
\end{contribuicao}

\begin{desenvolvimento}
\fichatecnica{
  cargo       = {},   % Cargo
  time        = {},   % Equipe liderada
  rituais     = {},   % Rituais e ferramentas (sprints, OKRs)
  documentos  = {},   % Documentos de autoria (códigos) — opcional
  indicadores = {},   % Indicadores acompanhados
}

\preencher{O planejamento (OKRs, sprints), o acompanhamento, os processos e os documentos que você criou ou mudou, com os códigos.}
\end{desenvolvimento}

\begin{resultados}
\preencher{Os indicadores: entregas no prazo, pessoas formadas, riscos tratados. O que ficou pronto para quem veio depois.}
\end{resultados}

\begin{evidencias}
% Uma linha por evidência: \evidencia{tipo}{identificador}{o que mostra}{onde conferir}
% \evidencia{Registro}{OKRs de <AAAA/S>}{Os objetivos do semestre e o fechamento}{SOMA › OKRs}
% \evidencia{Arquivo controlado}{NRO-PES-<SN> Rev. B}{O procedimento que você revisou}{SOMA › Arquivos}
\end{evidencias}

\begin{aprendizados}
\preencher{O que liderar ensinou, e que conteúdos (gestão de projetos, economia, administração) você aplicou.}
\end{aprendizados}
\end{atividade}
